\begin{tikzpicture}[scale=2]

% === Parâmetros ===
\def\r{3}   % raio do círculo
\def\a{2}   % raio do cilindro

% === Cálculos automáticos ===
\pgfmathsetmacro{\yint}{sqrt(\r*\r - \a*\a)} % coordenada y de interseção
\pgfmathtruncatemacro{\yintb}{\r*\r - \a*\a} % O que vai dentro da raiz, pra impressão
\pgfmathsetmacro{\ang}{atan(\yint/\a)}       % ângulo correspondente


  \begin{axis}[      
      xlabel={$y$}, ylabel={$z$},   
      axis equal image, % mantém escala x=y para círculos não ficarem ovalados
      every axis x label/.style={at={(ticklabel* cs:0.97)}, anchor=north},
      every axis y label/.style={at={(ticklabel* cs:0.97)}, anchor=east},
      xmin=-1.2*\r, xmax=1.2*\r, ymin=-1.2*\r, ymax=1.2*\r,
      % width=8cm, height=6cm,
      xtick={-\r,-\a,0,\a,\r},
      % xticklabels={-3,-2,0,2,3}
      ytick={-\r,-\yint,0,\yint,\r}, 
      yticklabels={-\r,$-\sqrt{\yintb}$,0,$\sqrt{\yintb}$,\r},
      x tick label style={anchor=north east}, 
      y tick label style={anchor=south east}, 
    ]

    % --- Curvas Parametrizadas ---

    \draw[thick, blue] (axis cs:0,0) circle [radius=\r];    % Círculo de raio r

    \addplot[thick, red, domain=-1.2*\r:1.2*\r] ({\a}, {x});      % Reta vertical x=a
    \addplot[thick, red, domain=-1.2*\r:1.2*\r] ({-\a}, {x});     % Reta vertical x=-a

    % --- Região 1 sombreada com Arc ---
    \path[fill=gray!30, opacity=0.6]
      (axis cs:\a,-\yint)     % Ponto inicial, 
      % necessário para o comando Arc a seguir.
      % Se vc adicionar um "--" depois do ponto acima, 
      % o "arc" não vai saber de onde começar
      arc[  start angle=-\ang, % \alpha = \arctan \left(\frac{\sqrt{5}}{2}\right)
            end angle=\ang,      % \betha = \pi - \alpha
            radius=\r]        % Arco do ciclo
      -- cycle;          % Conecta o ponto inicial ao final


    % --- Região 2 sombreada com Arc ---
    \path[fill=gray!30, opacity=0.6]
      (axis cs:-\a,\yint)     % Ponto inicial, 
      % necessário para o comando Arc a seguir.
      % Se vc adicionar um "--" depois do ponto acima, 
      % o "arc" não vai saber de onde começar
      arc[  start angle=(180-\ang), % \alpha = \arctan \left(\frac{\sqrt{5}}{2}\right)
            end angle=(180+\ang),      % \betha = \pi - \alpha
            radius=\r]        % Arco do ciclo
      -- cycle;          % Conecta o ponto inicial ao final

    % --- Linhas Tracejadas ---
     \addplot[Tracejado] coordinates {(-\a,{-\yint}) (\a,{-\yint})};
     \addplot[Tracejado] coordinates {(-\a,{\yint}) (\a,{\yint})};
     
    % --- Pontos de interseção ---
    \addplot[only marks, black!60, mark=*] coordinates {(-\a,{\yint})
                                                        (-\a,{-\yint})
                                                        (\a,{-\yint})
                                                        (\a,{\yint})
    };


  \end{axis}
  
\end{tikzpicture}
